\documentclass[12pt, titlepage]{article}

\usepackage{amsmath, mathtools}

\usepackage[round]{natbib}
\usepackage{amsfonts}
\usepackage{amssymb}
\usepackage{graphicx}
\usepackage{colortbl}
\usepackage{xr}
\usepackage{hyperref}
\usepackage{longtable}
\usepackage{xfrac}
\usepackage{tabularx}
\usepackage{float}
\usepackage{siunitx}
\usepackage{booktabs}
\usepackage{multirow}
\usepackage[section]{placeins}
\usepackage{caption}
\usepackage{fullpage}
\usepackage{amssymb}

\hypersetup{
bookmarks=true,     % show bookmarks bar?
colorlinks=true,       % false: boxed links; true: colored links
linkcolor=red,          % color of internal links (change box color with linkbordercolor)
citecolor=blue,      % color of links to bibliography
filecolor=magenta,  % color of file links
urlcolor=cyan          % color of external links
}

\usepackage{array}

\externaldocument{../../SRS/SRS}

\input{../../Comments.text}
\input{../../Common.text}

\begin{document}

\title{Module Interface Specification for \progname{}}

\author{\authname}

\date{\today}

\maketitle

\pagenumbering{roman}

\section{Revision History}

\begin{tabularx}{\textwidth}{p{3cm}p{2cm}X}
\toprule {\bf Date} & {\bf Version} & {\bf Notes}\\
\midrule
Date 1 & 1.0 & Notes\\
Date 2 & 1.1 & Notes\\
\bottomrule
\end{tabularx}

~\newpage

\section{Symbols, Abbreviations and Acronyms}

See SRS Documentation at \wss{give url}

\wss{Also add any additional symbols, abbreviations or acronyms}

\newpage

\tableofcontents

\newpage

\pagenumbering{arabic}

\section{Introduction}

The following document details the Module Interface Specifications for
\wss{Fill in your project name and description}

Complementary documents include the System Requirement Specifications
and Module Guide.  The full documentation and implementation can be
found at \url{...}.  \wss{provide the url for your repo}

\section{Notation}

\wss{You should describe your notation.  You can use what is below as
  a starting point.}

The structure of the MIS for modules comes from \citet{HoffmanAndStrooper1995},
with the addition that template modules have been adapted from
\cite{GhezziEtAl2003}.  The mathematical notation comes from Chapter 3 of
\citet{HoffmanAndStrooper1995}.  For instance, the symbol := is used for a
multiple assignment statement and conditional rules follow the form $(c_1
\Rightarrow r_1 | c_2 \Rightarrow r_2 | ... | c_n \Rightarrow r_n )$.

The following table summarizes the primitive data types used by \progname. 

\begin{center}
\renewcommand{\arraystretch}{1.2}
\noindent 
\begin{tabular}{l l p{7.5cm}} 
\toprule 
\textbf{Data Type} & \textbf{Notation} & \textbf{Description}\\ 
\midrule
character & char & a single symbol or digit\\
integer & $\mathbb{Z}$ & a number without a fractional component in (-$\infty$, $\infty$) \\
natural number & $\mathbb{N}$ & a number without a fractional component in [1, $\infty$) \\
real & $\mathbb{R}$ & any number in (-$\infty$, $\infty$)\\
\bottomrule
\end{tabular} 
\end{center}

\noindent
The specification of \progname \ uses some derived data types: sequences, strings, and
tuples. Sequences are lists filled with elements of the same data type. Strings
are sequences of characters. Tuples contain a list of values, potentially of
different types. In addition, \progname \ uses functions, which
are defined by the data types of their inputs and outputs. Local functions are
described by giving their type signature followed by their specification.

\section{Module Decomposition}

The following table is taken directly from the Module Guide document for this project.

\begin{table}[h!]
\centering
\begin{tabular}{p{0.3\textwidth} p{0.6\textwidth}}
\toprule
\textbf{Level 1} & \textbf{Level 2}\\
\midrule

{Hardware-Hiding} & ~ \\
\midrule

\multirow{7}{0.3\textwidth}{Behaviour-Hiding} & Input Parameters\\
& Output Format\\
& Output Verification\\
& Temperature ODEs\\
& Energy Equations\\ 
& Control Module\\
& Specification Parameters Module\\
\midrule

\multirow{3}{0.3\textwidth}{Software Decision} & {Sequence Data Structure}\\
& ODE Solver\\
& Plotting\\
\bottomrule

\end{tabular}
\caption{Module Hierarchy}
\label{TblMH}
\end{table}

\newpage
~\newpage

\section{MIS of Inference Engine} \label{mis:ie}

\subsection{Module}

Inference Engine

\subsection{Uses}

GGML (Section~\ref{mis:ggml}), Tensor I/O (Section~\ref{mis:tensorio})

\subsection{Syntax}

\subsubsection{Exported Types}

\begin{itemize}
  \item $\mathit{Report}$, as in the Core module.
  \item $\mathit{Tensor}$, as in the GGML module (Section~\ref{mis:ggml}).
  \item $\mathit{RequestID}$: $\mathbb{N}$. One inference request. A node keeps a
    separate key-value (KV) cache for each request.
  \item $\mathit{Step} = (r: \mathit{RequestID}, x: \mathit{Tensor}, \tau:
    \mathbb{R})$: one step of request $r$, received from Tensor I/O. $\tau$ is
    the request's sampling temperature.
\end{itemize}

\subsubsection{Exported Access Programs}

\begin{center}
\begin{tabular}{p{2.6cm} p{4.6cm} p{2.0cm} p{4.4cm}}
\toprule
\textbf{Name} & \textbf{In} & \textbf{Out} & \textbf{Exceptions} \\
\midrule
init & -- & -- & -- \\
load & $\mathit{file}$: string, $a: \mathbb{N}$, $b: \mathbb{N}$ & -- &
  InvalidRange, UnsupportedModel, LoadFailed, Running \\
unload & -- & -- & NotLoaded, Running \\
start & -- & -- & NotLoaded, Running \\
stop & -- & -- & -- \\
endRequest & $r: \mathit{RequestID}$ & -- & UnknownRequest \\
report & -- & $\mathit{Report}$ & NotInitialised \\
\bottomrule
\end{tabular}
\end{center}

{\raggedright \noindent The Core module calls every access program. Calling
load with a new range replaces the loaded partition. Between start and stop,
the engine repeatedly takes the next $\mathit{Step}$ from Tensor I/O, runs $x$
through units $a$ to $b$, and passes the result back to Tensor I/O. On the node
holding unit 1, $x$ is a sequence of token IDs. The node holding unit $N$
samples the next token at temperature $\tau$ and passes that on instead. A step
with the wrong shape, or one that would take its request past
\textit{CONTEXT\_MAX} tokens, is dropped along with that request's KV cache.
endRequest frees a finished request's KV cache.\par}

\subsection{Semantics}

\subsubsection{State Variables}

\begin{itemize}
  \item $\mathit{initialised}: \mathbb{B}$
  \item $\mathit{loaded}: \mathbb{B}$
  \item $\mathit{running}: \mathbb{B}$
  \item $\mathit{file}$: string, the GGUF file the partition is read from
  \item $a, b: \mathbb{N}$, the first and last partition units held
  \item $W$: string $\to \mathit{Tensor}$, the weights of units $a$ to $b$ by
    name
  \item $\mathit{kv}: \mathit{RequestID} \to \text{seq of } (k: \mathit{Tensor},
    v: \mathit{Tensor})$, the KV cache of each active request, one entry per
    transformer layer held
  \item $\mathit{count}: \mathbb{N}$, tokens processed since the last
    report
  \item $t_\mathit{last}: \mathbb{R}$, the time of the last report
\end{itemize}

\subsubsection{Environment Variables}

\begin{itemize}
  \item $t: \mathbb{R}$, the current time in seconds.
  \item The device's GPU-accessible memory, reached through the GGML module.
\end{itemize}

\subsubsection{Exceptions}

\begin{itemize}
  \item NotInitialised: any access program other than init is called before
    init. Every access program other than init can raise it.
  \item InvalidRange: $a > b$, or $b$ is greater than the number of partition
    units in the model.
  \item UnsupportedModel: the model's architecture is not one the Inference
    Engine can run.
  \item LoadFailed: the weights of units $a$ to $b$ do not fit in the node's
    GPU-accessible memory, or the file cannot be read.
  \item NotLoaded: start or unload is called when no partition is loaded.
  \item Running: load, unload or start is called between start and stop.
  \item UnknownRequest: endRequest names a request with no KV cache on this
    node.
\end{itemize}

\newpage

\section{MIS of GGML} \label{mis:ggml}

\subsection{Module}

GGML

\subsection{Uses}

None.

\subsection{Syntax}

\subsubsection{Exported Types}

\begin{itemize}
  \item $\mathit{Backend} = \{\mathsf{cpu}, \mathsf{metal}, \mathsf{cuda},
    \mathsf{vulkan}, \mathsf{rocm}\}$
  \item $\mathit{QuantType}$: the quantization types in the GGML version we
    use, from eight bits down to about one bit per weight (for example
    $\mathsf{q8\_0}$, $\mathsf{q4\_K}$, $\mathsf{iq4\_xs}$, $\mathsf{iq3\_s}$,
    $\mathsf{iq1\_s}$ and $\mathsf{nvfp4}$).
  \item $\mathit{DType} = \{\mathsf{f32}, \mathsf{f16}, \mathsf{bf16}\} \cup
    \mathit{QuantType}$: 32-bit or 16-bit floating point, or a quantization type.
  \item $\mathit{Tensor} = (\mathit{shape}: \text{seq of } \mathbb{N},
    \mathit{type}: \mathit{DType})$: an array held in the backend's memory.
  \item $\mathit{Op} = \{\mathsf{add}, \mathsf{mul}, \mathsf{matmul},
    \mathsf{norm}, \mathsf{rope}, \mathsf{softmax}, \mathsf{reshape}\}$
  \item $\mathit{Graph} = \text{seq of } (\mathit{op}: \mathit{Op},
    \mathit{args}: \text{seq of } \mathit{Tensor})$: operations in the order
    they run. An entry can use the results of earlier entries, and the last
    entry's result is the output.
\end{itemize}

\subsubsection{Exported Access Programs}

\begin{center}
\begin{tabular}{p{2.6cm} p{4.4cm} p{2.6cm} p{4.0cm}}
\toprule
\textbf{Name} & \textbf{In} & \textbf{Out} & \textbf{Exceptions} \\
\midrule
backends & -- & set of $\mathit{Backend}$ & -- \\
init & $\beta: \mathit{Backend}$ & -- & BackendUnavailable \\
readMetadata & $\mathit{file}$: string & string $\to$ string & FileError,
  InvalidFile \\
readTensors & $\mathit{file}$: string, $\mathit{names}$: set of string &
  string $\to \mathit{Tensor}$ & FileError, InvalidFile, UnknownTensor,
  OutOfMemory \\
alloc & $s$: seq of $\mathbb{N}$, $d: \mathit{DType}$ & $\mathit{Tensor}$ &
  OutOfMemory \\
compute & $g: \mathit{Graph}$ & $\mathit{Tensor}$ & InvalidGraph, OutOfMemory
  \\
release & $X$: set of $\mathit{Tensor}$ & -- & UnknownTensor \\
freeMemory & -- & $\mathbb{N}$ & -- \\
\bottomrule
\end{tabular}
\end{center}

{\raggedright \noindent Only the Inference Engine calls this module.
readMetadata returns a GGUF file's metadata, such as the architecture and
layer count. readTensors loads only the named tensors, so a node never reads
weights outside its partition.\par}

\subsection{Semantics}

\subsubsection{State Variables}

\begin{itemize}
  \item $\mathit{initialised}: \mathbb{B}$
  \item $\beta: \mathit{Backend}$, the backend in use
  \item $\mathit{held}$: set of $\mathit{Tensor}$, the tensors currently
    allocated
\end{itemize}

\subsubsection{Environment Variables}

\begin{itemize}
  \item The device's CPU and GPU memory, where tensors are stored.
  \item The file system, from which GGUF files are read.
\end{itemize}

\subsubsection{Exceptions}

\begin{itemize}
  \item NotInitialised: an access program other than backends or init is
    called before init.
  \item BackendUnavailable: $\beta$ is not available on the device.
  \item FileError: $\mathit{file}$ does not exist or cannot be read.
  \item InvalidFile: $\mathit{file}$ is not a valid GGUF file.
  \item UnknownTensor: a name in $\mathit{names}$ is not in the file, or a
    tensor in $X$ is not held.
  \item OutOfMemory: the new tensors do not fit in the backend's memory.
  \item InvalidGraph: an entry's arguments have the wrong shape or type for
    its operation.
\end{itemize}

\newpage

\section{MIS of Tensor I/O} \label{mis:tensorio}

\subsection{Module}

Tensor I/O

\subsection{Uses}

GGML (Section~\ref{mis:ggml})

\subsection{Syntax}

\subsubsection{Exported Types}

\begin{itemize}
	\item $\mathit{SocketAddrV4} = (
		\mathit{host}: \mathit{IPAddress},
		\mathit{port}: \{p \in \mathbb{N}_0|0 \leq p \leq 65535\}
		)$
	\item $\mathit{SocketAddrV6} = (
		\mathit{host}: \mathit{IPAddress},
		\mathit{port}: \{p \in \mathbb{N}_0|0 \leq p \leq 65535\},
		\mathit{scope}: \mathbb{N} \cup \mathit{null}
		)$
	\item $\mathit{SocketAddr} = \mathit{SocketAddrV4} \cup \mathit{SocketAddrV6}$
	\item $\mathit{Step} = (
		\mathit{requestID}: string,
		\mathit{activations}: \mathit{Tensor},
		\mathit{temperature}: \{t \in \mathbb{R}|0 \leq t \leq 2\}
		)$
\end{itemize}

\noindent The $\mathit{Tensor}$ type belongs to GGML.

\subsubsection{Exported Access Programs}

\begin{center}
\begin{tabular}{p{3cm} p{4.5cm} p{1.5cm} p{3.9cm}}
\toprule
\textbf{Name} & \textbf{In} & \textbf{Out} & \textbf{Exceptions} \\
\midrule
connect & $\mathit{upstream:} \mathit{SocketAddr}$, $\mathit{downstream:} \mathit{SocketAddr}$ & -- &
	\mbox{ConnectionRefused}, \mbox{Unreachable}, \mbox{ConnectionTimeout}, \mbox{ConnectionReset}, \mbox{AddressInUse},
	\mbox{AddressNotAvailable}, \mbox{ExistingConnection} \\
disconnect & -- & -- & \mbox{NoExistingConnection} \\
send & $\mathit{step}: \mathit{Step}$ & -- &
	\mbox{NoExistingConnection}, \mbox{ConnectionTimeout}, \mbox{ConnectionReset} \\
recv & -- & $\mathit{Step}$ &
	\mbox{NoExistingConnection}, \mbox{ConnectionTimeout}, \mbox{ConnectionReset} \\
\bottomrule
\end{tabular}
\end{center}

\noindent connect, disconnect, send, and recv are called by the Inference Engine.

\subsection{Semantics}

\subsubsection{State Variables}

\begin{itemize}
	\item $\mathit{connected}: \mathbb{B}$, whether this node is connected to upstream/downstream
		peers.
	\item $\mathit{upstream}: \mathit{SocketAddr} \cup \mathit{null}$, the upstream node's address.
	\item $\mathit{downstream}: \mathit{SocketAddr} \cup \mathit{null}$, the downstream node's address.
\end{itemize}

\subsubsection{Environment Variables}

\begin{itemize}
	\item The device's networking stack, through which upstream/downstream nodes will be connected.
\end{itemize}

\subsubsection{Exceptions}
\begin{itemize}
	\item ExistingConnection: Connections with upstream/downstream already exist. Must first
		disconnect before calling connect.
	\item NoExistingConnection: Connections with upstream/downstream don't exist. Must first connect
		before calling disconnect, send, or recv.
	\item ConnectionRefused: Peer is reachable, but nothing is listening on the given port.
	\item Unreachable: No route to given host.
	\item ConnectionTimeout: Response from peer is not received within allowed time.
	\item ConnectionReset: Peer closed the connection abruptly.
	\item AddressInUse: Attempted to bind to an endpoint is already bound by another application.
	\item AddressNotAvailable: Attempted to bind to an invalid endpoint.
\end{itemize}

\newpage

\bibliographystyle {plainnat}
\bibliography {../../../refs/References}

\newpage

\newpage{}

\section{Appendix --- Generative AI Use Disclosure}

\wss{ChatGPT (version 5) was used to brainstorm ideas for the template for this disclose section.}

\subsection{AI Tools Used}

\wss{Give the name of the AI tool(s) used and the version.}

\subsection{How and Where Used}

\wss{Explain what the tool(s) were used for.  Examples include: brainstorming,
explaining a technical concept, reviewing the document, generating or modifying
text, generating or modifying diagrams, identifying errors or omissions or
inconsistencies.}

\wss{For each of the uses, identify which parts of the document were affected.}

\wss{You do not need to report every interaction, but you should disclose the 
uses that materially contributed to your work}

\begin{itemize}
    \item \textbf{MIS of Inference Engine and MIS of GGML}: Anthropic Claude
    Opus 5.5, used through Claude Code, drafted both specifications from the
    SRS Inference Engine requirements, the Core module's interface and the Rust
    crates in the repository. We reviewed the draft and decided which module
    reads GGUF files and owns the tensor type.
\end{itemize}

\subsection{Verification of AI-Generated Content}

\wss{\begin{itemize}
    \item Describe how the team checked AI-generated or AI-assisted material
    before including it in the document.
    \item In particular, verify factual claims, technical assertions,
    requirements, calculations, references, and other information for which an
    AI system may produce incorrect or fabricated results.
    \item \textbf{The team is responsible for the accuracy, completeness,
consistency, and appropriateness of the submitted document, regardless of
whether material was generated by a human or an AI system.}
\end{itemize}}

We checked the Inference Engine's load and report against the calls the Core
module makes, and checked that the Inference Engine uses every GGML access
program.

\newpage{}

\section{Appendix --- Reflection}

\wss{Not required for CAS 741 projects}

The information in this section will be used to evaluate the team members on the
graduate attribute of Problem Analysis and Design.

\input{../../Reflection.text}

\begin{enumerate}
  \item What went well while writing this deliverable? 
  \item What pain points did you experience during this deliverable, and how
    did you resolve them?
  \item Which of your design decisions stemmed from speaking to your client(s)
  or a proxy (e.g. your peers, stakeholders, potential users)? For those that
  were not, why, and where did they come from?
  \item While creating the design doc, what parts of your other documents (e.g.
  requirements, hazard analysis, etc), it any, needed to be changed, and why?
  \item What are the limitations of your solution?  Put another way, given
  unlimited resources, what could you do to make the project better? (LO\_ProbSolutions)
  \item Give a brief overview of other design solutions you considered.  What
  are the benefits and tradeoffs of those other designs compared with the chosen
  design?  From all the potential options, why did you select the documented design?
  (LO\_Explores)
  \item (After you have implemented another team's module, which means this
  isn't filled in until after the original deadline). What did you learn by
  implementing another team's module? Were all the details you needed in the
  documentation, or did you need to make assumptions, or ask the other team
  questions?  If your team also had another team implement one of your modules,
  what was this experience like?  Are there things in your documentation you
  could have changed to make the process go more smoothly for when an
  ``outsider'' completes some of the implementation?

  Your reflection should be about concepts learned in your SE education, such as modular decomposition, cohesion, coupling, information hiding, etc.  You should critique the quality and utility of the specification you received.

  You should create a PR for your implemented module.  A link to the PR should be part of your reflection.

\end{enumerate}

\end{document}
